\documentclass[../../../include/open-logic-section]{subfiles}

\begin{document}
	
\olfileid{sth}{replacement}{finiteaxiomatizability}
\olsection{ملحق: قابلية الوضع في مسلّمات منتهية}
نختم الفصل بنتائج تستخرج من الاستبدال. أولها لـ\citet{Montague1961}؛
وينبغي التنبه إلى أن البرهان ليس \emph{داخل}~$\ZF$، بل هو برهان
\emph{عن}~$\ZF$:
\begin{thm}\ollabel{zfnotfinitely}
	لا يمكن وضع~$\ZF$ في مسلّمات منتهية. وأعم من ذلك: إذا كانت~$\Th{T}$
	منتهية وكان~$\Th{T} \Proves \ZF$، كانت~$\Th{T}$ غير متسقة.
	
	(ونقصر النظر ضمنًا على جمل الرتبة الأولى التي ليس في رمزها الأولي
	من غير المنطقيات إلا~$\in$، ونقصد بـ$\Th{T} \Proves \ZF$ أن
	$\Th{T} \Proves \phi$ لكل~$\phi \in \ZF$.)
\end{thm}
\begin{proof}
	ثبت نظرية منتهية~$\Th{T}$ بحيث~$\Th{T} \Proves \ZF$. ومن ثم تثبت
	$\Th{T}$ الانعكاس، أعني~\olref[sth][replacement][ref]{reflectionschema}.
	ولما كانت~$\Th{T}$ منتهية، جاز جمعها في اقتران واحد~$\theta$.
	وبعكس هذه الصيغة ينتج
	$\Th{T} \Proves \exists \beta(\theta \liff \theta^{{\OLId{latin-looped}{V}}_\beta})$.
	ومن البين أن~$\Th{T} \Proves \theta$، ومن ثم
	$\Th{T} \Proves \exists \beta\ \theta^{{\OLId{latin-looped}{V}}_\beta}$.
	
	ولنرمز بـ$\psi({\OLId{latin-looped}{X}})$ اختصارًا إلى:
	\[
		\theta^{\OLId{latin-looped}{X}} \land {\OLId{latin-looped}{X}}\text{ مرحلة} \land (\forall {\OLId{latin-looped}{Y}} \in {\OLId{latin-looped}{X}})({\OLId{latin-looped}{Y}}\text{ مرحلة}\lif \lnot \theta^{{\OLId{latin-looped}{Y}}})
	\]
	ومعناه التقريبي أن~${\OLId{latin-looped}{X}}$ نموذج مرحلي لـ$\theta$، وهو أدنى من هذه
	الجهة بحسب~$\in$. وإذا استحضرنا
	$\Th{T} \Proves \exists \beta\ \theta^{{\OLId{latin-looped}{V}}_\beta}$، واستعملنا خواص
	الرتب الأساسية داخل~$\ZF$، ومن ثم داخل~$\Th{T}$، حصلنا على:
	\begin{equation}
		\Th{T} \Proves \exists {\OLId{latin-looped}{M}} \psi({\OLId{latin-looped}{M}}). \tag{*}\label{Mpsi}
	\end{equation}
	وباستخدام المقترن الأول من~$\psi({\OLId{latin-looped}{X}})$، متى كان~$\Th{T} \Proves \sigma$،
	كان~$\Th{T} \Proves \forall {\OLId{latin-looped}{X}}(\psi({\OLId{latin-looped}{X}}) \lif \sigma^{\OLId{latin-looped}{X}})$.
	وبمقتضى~\eqref{Mpsi}:
	\begin{align*}
		\Th{T} &\Proves \forall {\OLId{latin-looped}{X}}(\psi({\OLId{latin-looped}{X}}) \lif (\exists {\OLId{latin-looped}{N}} \psi({\OLId{latin-looped}{N}}))^{\OLId{latin-looped}{X}})\\
	\intertext{وباستخدام هذا، ثم \eqref{Mpsi} مرة أخرى:}
		\Th{T} &\Proves \exists {\OLId{latin-looped}{M}}(\psi({\OLId{latin-looped}{M}}) \land (\exists {\OLId{latin-looped}{N}} \psi({\OLId{latin-looped}{N}}))^{\OLId{latin-looped}{M}})
	\intertext{ومن ثم، بوجه خاص:}
		\Th{T} &\Proves \exists {\OLId{latin-looped}{M}}(\psi({\OLId{latin-looped}{M}}) \land (\exists {\OLId{latin-looped}{N}} \in {\OLId{latin-looped}{M}})(({\OLId{latin-looped}{N}}\text{ مرحلة})^{\OLId{latin-looped}{M}} \land (\theta^{\OLId{latin-looped}{N}})^{\OLId{latin-looped}{M}}))
	\intertext{وباستدلال أولي عن إطلاق المراحل والرتب:}
		\Th{T} &\Proves \exists {\OLId{latin-looped}{M}}(\psi({\OLId{latin-looped}{M}}) \land (\exists {\OLId{latin-looped}{N}} \in {\OLId{latin-looped}{M}})({\OLId{latin-looped}{N}}\text{ مرحلة} \land \theta^{\OLId{latin-looped}{N}}))
	\end{align*}
	فتكون~$\Th{T}$ غير متسقة.
	\footnote{يتطلب هذا «الاستدلال الأولي»
	إثبات بعض حقائق الإطلاق المتعلقة بمراحل الهرم ورتبه.}
\end{proof}

وهذه نتيجة من الجنس نفسه، نبه إليها~\citet[223]{Potter2004}:

\begin{prop}\ollabel{finiteextensionofZ}
	لتكن~$\Th{T}$ امتدادًا لـ$\Z$ بعدد منته من المسلّمات الجديدة. إذا
	كان~$\Th{T} \Proves \ZF$، فإن~$\Th{T}$ غير متسقة. (والقيود
	الضمنية هنا هي نفسها التي في~\olref{zfnotfinitely}.)
\end{prop}
\begin{proof}
	اجعل~$\theta$ اقتران مسلّمات~$\Th{T}$ كلها \emph{إلا} مثيلات الفصل
	غير المتناهية. وعرّف~$\psi$ من~$\theta$ كما في
	\olref{zfnotfinitely}. فعندئذ يمكننا، استنادًا إلى المراحل، إثبات
	$\Th{T} \Proves \exists {\OLId{latin-looped}{M}} \psi({\OLId{latin-looped}{M}})$.
	
	وكما في~\olref{zfnotfinitely}، نستطيع إثبات المخطط الآتي: إذا كان
	$\Th{T} \Proves \sigma$، كان
	$\Th{T} \Proves \forall {\OLId{latin-looped}{X}}(\psi({\OLId{latin-looped}{X}}) \lif \sigma^{\OLId{latin-looped}{X}})$. وبعد ذلك يتم
	البرهان تمامًا كما في~\olref{zfnotfinitely}.
	
	غير أن إقامة هذا المخطط تحتاج إلى زيادة يسيرة في العمل عما احتاج
	إليه في~\olref{zfnotfinitely}. فمثيلات
	الفصل داخلة في~$\Th{T}$، وليست من مقترنات~$\theta$. وندفع هذا بأن
	نثبت أن~$\Th{T} \Proves \forall {\OLId{latin-looped}{X}}({\OLId{latin-looped}{X}}\text{ مرحلة} \lif \sigma^{\OLId{latin-looped}{X}})$
	لكل مثيل فصل~$\sigma$؛ ونكل تفصيله إلى القارئ.
\end{proof}
\begin{prob}
	بيّن أنه لكل مثيل فصل~$\sigma$ يكون:
	$\Z \Proves \forall {\OLId{latin-looped}{X}}({\OLId{latin-looped}{X}}\text{ مرحلة} \lif \sigma^{\OLId{latin-looped}{X}})$.
	(استخدمنا هذا المخطط في
	\olref[sth][replacement][finiteaxiomatizability]{finiteextensionofZ}.)
\end{prob}
\begin{prob}
	بيّن أنه لكل~$\phi \in \Z$ يكون
	$\ZF \Proves \phi^{{\OLId{latin-looped}{V}}_{\omega+\omega}}$.
\end{prob}
\begin{prob}
	تحقق من سائر النتائج المخططية التي استعملت في برهاني
	\olref[sth][replacement][finiteaxiomatizability]{zfnotfinitely} و
	\olref[sth][replacement][finiteaxiomatizability]{finiteextensionofZ}.
\end{prob}

وقد أشرنا في~\olref[sth][replacement][absinf]{sec} إلى أن هذا يثبت
أن الاستبدال أقوى تمامًا من
\olref[ordinals][ordtype]{thmOrdinalRepresentation}. وبتعبير أدق:
إذا كانت~$\Z$ + «كل ترتيب جيد متماثل مع عدد ترتيبي فريد» متسقة، لم
تقدر على إثبات بعض مثيلات الاستبدال.


	%	By assumption, $\psi^{V_\alpha}$ has the form:
	%	\[
	%		(\forall A \in V_\alpha)(\exists S \in V_\alpha)(\forall x \in V_\alpha)(x \in S \liff (\phi^{V_\alpha}(x) \land x \in A))
	%	\]
	%	To establish this holds, fix $A \in V_\alpha$. Using Separation, obtain:
	%	\[
	%		S = \Setabs{x \in A}{\phi^{V_\alpha}(x)}
	%	\]
	%	Now $S \in V_\alpha$, since $S \subseteq A \in V_\alpha$, and clearly $(\forall x \in V_\alpha)(x \in S \liff (\phi^{V_\alpha}(x) \land x \in A)$.

\end{document}
